% !TeX root = ../papers/nc-p3.tex
\section{片方の行列式二次式が恒等的に消える全境界}
\label{sec:single-zero-pairings}
\begin{theorem}\label{thm:single-zero-pairings}
定理\ref{thm:klein-all-pairings}の可逆性を保つチャートでは、
$P,Q$の少なくとも一方が零多項式なら標準AS正則点はない。
片方だけが零となる点は
$\lambda\in\{1,-1,1/z,-1/z\}$に限られ、その始点0の
長さ7の経路次元は少なくとも128である。標準値は120である。
\end{theorem}
\begin{proof}
$P\equiv0$の場合には定数項から$d=za$または$d=a/z$である。
$a=0$なら$d=f_0=g_0=0$で一次係数$m=1$ゆえ矛盾する。
従って$a=1$に規格化できる。$Q\equiv0$の場合には
$c=b$または$c=z^2b$、同様に$k=-1$となる場合を除いて$b=1$とできる。
それぞれ次の二つの指定値を使う。$T=Y_\lambda A_0^{-1} X_\lambda$とする。
\[
\begin{array}{c|c|c}
 &\text{第一の場合}&\text{第二の場合}\\\hline
 P\equiv0:\ (T_{00},T_{11}) &(1,1)/f_0&(1,z^2)/g_0\\
 Q\equiv0:\ (T_{01},T_{10}) &(1,z)/e_0&(z,1)/e_0
\end{array}
\]
これは零にする最高次係数と一次係数を解いたもので、分母は非零である。
$D=ad-bc$とし、取り出した二成分の分子を$N_1,N_2$、
表の指定値を$u_1,u_2$と書く。
$E=N_1-u_1D$, $G=N_2-u_2D$を同時に零とする。

まず四つの例外倍率を除く。$u_2E-u_1G$は未定二変数
$(b,c)$または$(a,d)$の一次式で、後ろの変数の係数は単元である。
それを消去して$E$を取ると、残る変数の二つの一次式の積に単元を掛けた式になる。
例えば$P\equiv0,d=z,f_1=0$のとき、$c$を解いた後の式は
\[
 \frac{(b-\lambda)(\lambda b-1)(\lambda z-1)(\lambda z+1)}
 {\lambda(\lambda-1)(\lambda+1)}
\]
である。どの根でも他方の行列式二次式の三係数が全て消える。
\path{computations/verify_single_zero_pairings.py}は二種類の零多項式、
二つの比、表の二選択の各場合に消去を行い、
除数と二次式の最高次係数が局所化した環の単元であること、および両根での
三係数の消滅を検算する。従ってこれらは前節で全て排除済みである。

例外倍率を一つ固定すると、$u_2E-u_1G$は二変数の一方だけの一次式になる。
その係数の分子・分母の素因子は$z(z^2-1)$の因子に限られ、
この変数を解くと$E,G$はともに恒等的に零になる。残る変数を$u$とする。
例えば上の$P$の例で$\lambda=1$なら$b=1$、$u=c$である。
$Q$が非零で解を持つなら、判別式が零なので$Q$は真の二次式で、
$y=-k/(2q)$である。非零定数では解がなく、一次式なら判別式は非零だからである。
$x=v$は自由であり、四つの残りのSchur等式は恒等式になる。
$Q\equiv0$なら$x=-m/(2p)$, $y=v$として同じ議論を使う。
これで片方だけ零の全許容点は、二種類・二つの比・二選択・四倍率の
計32表示の各々で$(z,u,v)$のアフィン空間の開部分に入る。
この表示では重根の分母が零になる点を単に連続性で埋めず、
もう一式も零になる場合を前節、零にならない場合を上の次数比較で処理した。

同コードで各表示の$\mathbb Q(z,u,v)$上の四窓を作り、始点0で
長さ7までの語商を計算する。各表示で長さ0から6の次元は順に
$1,4,10,20,35,56,84$である。長さ7では128または134となる。
六関係の基底選択に使う点は$(z,u,v)=(3,5,7)$であり、
関数体上の階数計算はこの点の数値を用いない。
既約な表示上で特殊化した次元はこの値以上であるから、
各点で$\binom{10}{3}=120$を超える。
\end{proof}
